% !TeX spellcheck = en_US
\documentclass[french]{yLectureNote}

\title{Introduction à la termodynamique}
\subtitle{Thermodynamique}
\author{Paulhenry Saux}
\date{\today}
\yLanguage{Français}

\professor{}

\usepackage{graphicx}%----pour mettre des images
\usepackage[utf8]{inputenc}%---encodage
\usepackage{geometry}%---pour modifier les tailles et mettre a4paper
%\usepackage{awesomebox}%---pour les boites d'exercices, de pbq et de croquis ---d\'esactiv\'e pour les TP de PC
\usepackage{tikz}%---pour deiffner + d\'ependance de chemfig
\usepackage{tkz-tab}
\usepackage{tabularx}%---pour dimensionner automatiquement les tableaux avec variable X
\usepackage{awesomebox}%---Pour les boites info, danger et autres
\usepackage{menukeys}%---Pour deiffner les touches de Calculatrice
\usepackage{fancyhdr}%---pour les en-t\^ete personnalis\'ees
\usepackage{blindtext}%---pour les liens
\usepackage{hyperref}%---pour les liens (\`a mettre en dernier)
\usepackage{caption}%---pour la francisation de la l\'egende table vers Tableau
\usepackage{pifont}
\usepackage{array}%---pour les tableaux
\usepackage{lipsum}
\usepackage{yFlatTable}
\usepackage{multicol}
\newcommand{\Lim}[1]{\lim\limits_{\substack{#1}}\:}
\renewcommand{\vec}{\overrightarrow}
\newcommand{\bz}{\overline{z}}
\DeclareMathOperator{\card}{card}
\renewcommand{\vec}{\overrightarrow}
\newcommand{\N}[0]{\mathbb{N}}
\newcommand{\R}[0]{\mathbb{R}}
\newcommand{\C}[0]{\mathbb{C}}
\newcommand{\dd}[0]{\mathrm{d}}
\newcommand{\er}{\vec{e_{\rho}}}
\newcommand{\ep}{\vec{e_{\varphi}}}
\newcommand{\pip}{\pi^+}
\newcommand{\pim}{\pi^-}
\newcommand{\mv}{\mathcal{V}}
\DeclareMathOperator{\Div}{Div}
\newcommand{\rot}{\vec{\mathrm{rot}}}
\newcommand{\grad}{\vec{\mathrm{grad}}}
\begin{document}
\setcounter{chapter}{4}

\chapter{Approche  Thermodynamique de la transition de phase}
\section{Introduction}
\subsection{Diversité des transitions de phase}
\begin{definition}[Phase (général)]
    Pas un synonyme d'état de la matière. Un ensemble d'états accessibles à un corps, obéissant au meme jeu de loi de comportement. Synonyme : Phénoménologie. Un corps homogène n'a qu'une phase.
\end{definition}
\begin{definition}[Transition de phase]
    Changement de Phénoménologie au dela d'une valeu seuil d'un paramètre que l'on fixe.
\end{definition}
On considère qu'il n'existe qu'une sorte de  phase liquide et solide et gazeuse

\criticalInfo{Nomenclature des transitions de phase (par coeur)}{Noms à noter}

\begin{itemize}
    \item Changement d'état Solide-Liquide-Gaz
    \item Supraconductivité : On passe d'un matérieux avec une résistavité elctrique à un matérieux avec une résistivité nulle.
    \item Morphogénèse : appartition de forme d'états non uniformes
    \item Cristaux liquides : 
    \item Superfluides : un fluide refroidi proche de 0K, coule sans effort.
    \item Transition ferro-paramagnétiques : des solides chauffés perdent leur propriété d'aimantation
\end{itemize}
\subsection{2 expériences}
\subsubsection{Expérience avec volume imposé}
On réduit le volume avec pas constant (compression) au contat d'un thermostat à $T_0$. On attend l'équilibre pour chaque valeur de V. Variables fixées à l'équilibre : T et V. On mesure la pression p et on relève le volume.

Variante : On augmente la pression et on relève le volume.

Photo avec une image simplifiée de l'expérience : frontière V/VL/L et pression stable en VL

Courbe à faire p en fonction de V 
|
|
|_________ 
            -
L     L+V   V-
              -
              Il y a 3 Phénoménologie et 2 seuils en volume

              V en fonction de p 
-
  -


    ----------
Indiquer le pas.
Il existe ici qu'un seuil avec  2 Phénoménologies.


\subsection{Problématiques}
Pas de modèle adéquat pour traiter ces problèmes. Le modèle du GP ne convient pas : on ne préduti que la partie liquide qui ressemble à l'isotherme. Le liquide incompressible ne fonctionne pas non plus.

On est bloqué.

Et le modèle de vanderValls : (p + a/V^2 n^2)(V-nb) = nRT. Cela fait le lien entre les courbes mais on ne prédit pas les états diphasiques.

Les états au milieu de la courbe ne sont pas des états d'équilibre.

Comment introduire la tranisiton de phase dans le modèle ?

\begin{proposition}
    La tranisiton de phase sera possible des que l'état d'éq devendra moins stable qu'un autre état d'équilibre eccessible dans les mêmes conditions. Elle se fait assurément si l'état d'équilibre devient instable.
\end{proposition}
On utilisera les potentiels Thermodynamiques
\section{Stabilité de l'état d'équilibre}
\subsection{Définition qualitative}
Qualifier la stabilité de l'état d'équilibre c'est énoncer le comportement du système suite à une perturbation.

Il y a 2 cas : Soit l'état A est le seul accessible dans les conditions p,T (si on bouge, on revient vers A, au sens de LYAPUNOV) ou alors il y a  états A et B accessibles.

2e possibilité : il y a une évolution sptannée où l'on s'éloigne de A. A est alors un état d'équilibre instable.

3e possibilité : Si je perturbe avec une amplitude plus grande mais toujours finie. Si $\Delta V_p >\Delta V_{p,c},$ il y a évolution spontannée, dans le cas contraire, il y a retour. Cela dépend de la variation. A est alors un état meta-stable.

\begin{definition}[Stabilité]
    Retour vers A quelque soit l'amplitude de la perturbation. (quand il n'y a qu'un état accessible)
\end{definition}
En mécanique, l'état métastable sera une cuvette finie.

Quand il y a évolution spontannée, en métastable, on parle de rupture de métastabilité.

Une façon de quitter la stabilité est de faire tendre $\Delta V\to 0$.
\begin{definition}[Stabilité]
    A est stable si l'évolution spontannée depuis est A est impossible.
\end{definition}
En résumé : Instable si une pertutbation d'amplitude infinitésimale suffit à faire dégager
Metastable : Spontanné vers un autre état ou équilibre dans les memes conditions
\subsection{Approche formelle}
\subsubsection{Fonction potentiel thermodynamique}
\begin{definition}
    Fonction d'état telle que pour un jeu de contrainte donné, ses variations permettent de disucuter des des possibilités d'évolution spontannée vers un état d'équilibre.
    Elle prend une valeur minimale dans un état stable
\end{definition}
On va chercher ces fonctions pour des états d'équilibre à (T et p) et (T et V).
Pour l'univers, $\Delta S_v = S^p_{I\to F}$. Dans un état non contraint, l'état d'équilibre est donné pour le max de S. L'état s'équilibre est stable.

Le potentiel est alors $\Psi = -S$.

On cherche pour l'autre  partie : On libère un piston avec un gaz compressé.

On sait que $\Delta U = W+Q, W = -p\Delta V$

On utilise le second principe : $S = \Delta S+ \Delta S^{ext}= \Delta S_{if} - \frac{Q}{T_0}$ car le thermostat est une source.

On en déduit $S = \frac{\Delta U + p_0 \Delta V}{T_0}$

Soit la fonction d'état $G^\star = U+p_0 V - T_0 S$ donc $\Delta G^\star = -T_0S_{I\to F}^P = \Delta G$. avec $G = U+pV-TS$

L'évolution impossible serait que les états B aient un G plus grand que A. G_A est min pour l'état stable. G est l'éngier libre de Gipps.

\warningInfo{Différence entre G et G*}{G* : on utilise les valeurs de l'atm et G à l'équilibre}

On peut écrire $\Delta G = -T_0 S^p_{i\to f}$ sur les transformations spontannées

La transformation 2.b est monobare au contact d'un thermostat.

Au cours de l'évolution, G ne peut que décroitre, Si j'atteins un état d'équilibre tel que G est minimum, à partir de cet états, tous les $\Delta G$ seront positifs. Mais c'est interdit. Donc un état où $G$ est minium est un état où l'évolution spontanné est impossible. 

Or, par définition, un état d'équilibre à partir duquel l'évolution spontannée est impossible est un état d'EQ Stable.

Donc $G$ correspond à un potentiel des évolutions spontannées monobares au contact d'un thermostat.


\subsubsection{Enthalpie libre ou énergie libre de Gibas}
% Il existe une fonction potentiel associée à une évolution au contat d'une entmosphère dont la température et la pression snt fixées.
% \paragraph{Contact avec une atmosphère} On a une déroissance de G*.
% \paragraph{ENtre 2 EQ dans les memes containtes}

\subsubsection{Énergie libre F}
Si on regarde les évoutions spontannées isochores au contact d'un thermostat, on obtient que la production d'netropie vaut $F = U-TS$. Quand j'ai un état d'équilibre à 

\subsubsection{Critère de Gibbs-Duhem}
\begin{theorem}
 Si   $\Delta U+p\Delta V-T\Delta S \geq 0$ , l'état d'équilibre est stable
\end{theorem}
On ne l'appliqyera que avec $\Delta G\geq 0, \Delta F\geq 0$.

G est l'énergie libre de Gibs ou enthalpie libre et F est l'énergie libre ou énergie libre de Helmotz

\criticalInfo{title}{Savoir refaire la démonstration avce monobare et icochore au contact d'un thermostat}

\subsubsection{Conclusion}
% Penser a uploader le code pour le projet
\chapter{Transition de phase : Mise en oeuvre et application}
\section{
    Une première mise en oeuvre
}
\subsection{Expression de G en fonction des conditions de pression pour un T0 donné.}
Au départ, 
\begin{flalign}
    G &= U+pV-TS\\
    \dd G &= \dd U + p\dd V+V\dd p - T\dd S-S\dd T\\
    &\dd U &= \dd U_{qs} = \delta W_{qs} + \delta Q_{qs} = -p\dd V+T \dd S\\
    \dd G &= V\dd p - S\dd T 
\end{flalign}
avec V et S dépendant de p et T.

On en déduit $$\dd G_{T_0} = V\dd p$$ donc $$G_{T_0} = \int V\dd p + cst$$

Pour un GP, $$V = \frac{mrT_0}{p} \Rightarrow G_{T_0} = mrT_0 \ln(p)+cst$$


Maintenant pour un liquide incompressible, V = CST, donc $G_{T_0} = V_Lp+cst$ 

On peut tracer les courbes obtenues

On définie une valeur de pression seuil à la croisée des courbes. À gauche, l'évolution vers Liquide est possible, à l'inverse à droite, l'évolution vers , donc à droite vapeur stable et à droite liquide stable.

On peut tracer ka nouvelle courbe.



\subsection{Application du critère de Gibbs Duhem}

\subsection{Pression des saturation Etats saturants - Isothermes critiques}
\begin{definition}[Pression de saturation]
    Pour une isotherme, c'est la valeur seuil pour laquele on a le changement de stabilité est appelé. 
\end{definition}
\begin{definition}[Volume saturant]
    2 états identiquement stables peuvent exister : l'un liquide est l'autre stable, avec des volumes de stautration différents.
\end{definition}
La pression de saturation et les volumes sturants sont des fonctions de la température : $p_s(T)$ et $V_{sat}(T)$

À la température critique, est telle que $v_sat^G = v_sat^L$.

On passe alors continumenent entre les 2 états.

\section{Résutats supplémentaires avec une phase solide}
Avec une phase solide, 
on obtient de nouveaux graphes :

Pour la première courbe, on rajoute juste le solide et pour la deuxième il n'y a pas de phase liquide.

\section{Synthèse des résuktats de l'analyse de stabilité pour le corps homogène}
Ces ésultats sont très imporants.
\subsection{Diagramme P et T}
On note la courbe de vaporisation, de sublimation et de fusion.

La courbe de vaporisation est bornée. On parle alors de fluide supercritique.

%Tracer courbe à partir de point triple et point critique
Le résultat btenu avec les poten,tiels sont dans le diagramme  PT
\subsection{Relation de Clapeyron}
%------------ Cours du 08/04/2026----------------
\subsubsection{Chaleurs Latentes - Relation de Clapeyron}
Pour tout point d'une courbe se saturation, ici de vaporisation
\begin{definition}[Chaleur latente]
    $$l_v(T) = h_v(T)^{sat}-h_L^{sat}$$ aussi nommée enthalpie de vaporisation
\end{definition}
On peut la trouver sous forme massique ou molaire.
\begin{theorem}[Relation de Clapeyron]
    $$l_v(T) = T(v_V^{sat}-v_L^{sat})\frac{\dd p_{sat}}{\dd T}$$
\end{theorem}
La dérivée de la fonction $p_{sat}$, c'est la pente de la courbe de vaporisation. À partir du diagramme, j'obtiens une grandeur énergétique. Attention, ce sont des volumes massique ou molaires et non de simples volumes.

On obtient le m\^eme type de relations avec les autres courbes de saturation.

\subsubsection{Conséquence 1 }
Le signe de la chaleur latente et variation avec la température
\subsubsection{Conséquence 2}
Entre 2 états naturels, on a $ H_v-H_L$ et donc si on connait $l_v(T)$ alors on connait $(k_H)_v-(k_H)_L$ alors on sait calculer $H_v-H_l$ entre 2 états d'équilibre quelconques alors on sait calculer 
$u_v-u_l$ entre 2 états quelconques puisque $U=H-pV$

La conaissance de la chaleur latente permet de faire le lien entre les phénomélogies pour l'énergie.

Exemle.

On se place sur une vaporisation monobare. : $Q = \Delta H = H_V-H_L$. Avec les  phénoménologies, $H_V = (c_p)_vT+K_v$, $H_L = c_LT+p_{atm}V_L+K_L$

On trace $H=f(T)$ et on obtient $\Delta H = H_L(T_S)-H_i+ml_v(T_s)+H_F-H_v$.
\chapter{Transition de phase - États non homogènes}
\section{Description de l'état d'équilibre d'un corps où les états de matière coexistent}
Modèle diphasique = 2 corps homogènes séparés par une interface sans epaisseur et $V=Vl+V_v, m=m_l+m_v$

Par contre, $U \neq U_l+U_v$, il faut rajouter un terme : $U = U_L+U_v+U_i$ avec $U_i$ l'énergie de l'interface. Cependant, on fera l'apprxmation que $U_i$ est très petite, donc négligeable deavnt l'néergie mise en jue.

Cela vaut pour toutes les grandeur  $Y = Y_L+Y_m$.
\subsection{modèle d'un état d'équilibre diphasique}

\subsection{Nouvelle variable : le titre massique}
$m_v$ et $m_l$ ne sont pas 2 variables indépendantes, on a $m_L=m_m_v$.
\begin{definition}[titre massique]
    $$x = \frac{m_V}{m}$$
\end{definition}
\section{Étude de la stabiité de l'équilibre du corps diphasique}
\subsection{États d"équilibres : pression et volumes accessibles sur l'isotherme}
Quelles sont les conditions de pression pour lesquelles l'état diphasique peut-il exister ? pour une isotherme ?

On calcule G : $G_{T_0} = (1-x)(G_{T_0})_L+x(G_{T_0})_V$. Seule valeur de p telle que $G_{LV}$ test min, c'est $p_s$.

Donc $Y_{LV} = m(1-x)y_L^{sat}(T)+mxy_V^{sat} = (1-x)Y^{sat}_{tot,L} + x Y_{tot,v} = f(T,x)$

\subsubsection{Quelles valeurs de volume sont accessible à l'état diphasique}
Donc $V_l{sat}<V_{LV}<V_V^{sat}$

On peut tracer les isothermes d'Andrew

|
|
|_____________
                *  
                    *
                        *   
                            *
L'état diphasique n'existe à l'équilibre stable que dans les conditions de saturation.
Sur la courbe de droite, on est vapeur, sur celle de gauche, on est liquide.
\subsection{Stabilité }

\section{Synthèse des réuslyats de l'analyse}
\subsection{Équilibre liquide-vapeur}
\subsection{Diagramme avec phase solide}

\subsection{Important - Expression des fonctions d'état entre 2 états d'équilibre diphasiques}
\subsubsection{Rappels expressions des variables des fonctions d'état entre 2 états saturants homogènes}
\begin{flalign*}
    \Delta H_{12} &= ml_v(T_0)\\
    \Delta U &= \Delta H - \Delta(pV)\\
                ml_v(T_0)-p_{sat}(T_0)(V_v-V_L)\\
    \Delta S_{12} => H_V-T_0S_V = H_L-T_0S_L => \Delta S = \frac{H_V-H_L}{T_0} = \frac{ml_v(T_0)}{T_0}\\           
\end{flalign*}
Au dessus = À savoir par coeur
\subsubsection{Entre 2 états diphasiques sur le meme palier}
On écrit :
\begin{flalign*}
    H_1 &= mx_1h_v(T_0)+m(1-x_1)h_L(T_0)\\
    H_2 &= mx_2h_v(T_0)+m(1-x_2)h_L(T_0)\\
    H_2-H_1 = m(x_2-x_1)l_v(T_0)
    \Delta U &= \Delta x (U_V-U_L)\\
    &= \Delta x(ml_v(t_0)-p_{sat}(V_v-V_L))\\
    \Delta S &= \Delta x \frac{ml_v(T_0)}{T_0}
\end{flalign*}
\subsubsection{2 états quelconques}
On prend le chemin que l'on souhaite. Le but est d'avoir un chemin par pallier. 
1->1' : Liquéfaction sur isotherme T0 -> 3 : Transformation isotherme isochore d'un liquide ->2' : Isobare d'une liquide -> 2 : Vaporisation isotherme à T1
\begin{flalign*}
    \delta H_{11'}  &=-x_1ml_v\\
    \Delta H_{1'3} &=v_l(p_{sat}(T_1)-p_{sat}(T_0))\\
    \Delta H_{32'} &=  mc_l(T-1-T_0)\\
    \Delta H_{2'2} = mx_2l_v(T_1)\\
    \Delta H &= mx_2l_v(T_1)-mx_1l_v(T_0)+mc_l(T_1-T_0)+V_L(T_0)(p_{sat}(T_1)-p_{sat}(T_0))\\
    \Delta H &= mx_2l_v(T_1)-mx_1l_v(T_0)+\Delta H_L
\end{flalign*}

\subsection{Détermination du titre massique vapeur}
Rappel : $$Y_{LV} = m_Lxy_V^{sat}(T)+m(1-x)y_L^{sat} = f(T,x)$$
Si T est connue et $Y_{LV}$ est connue dans 1 état d'équilibre, aors x est connu.

Exemple avec $T_0, V_0$ connus.

Passez du titre au volume et vice versa est à savoir faire
\begin{theorem}[Théorème des moments]
    $$x_0 = \frac{LM}{LV}$$
\end{theorem}
\section{Exemples}
On prend $V_0 = 2ml (L+V)$ Trouver $x_0$
On prend $T=18\circ$ avec les données de la colonne à mercure :
On peut en déduire, une fois $x_0$ déterminé que $v_v = m_vv_v = xV_v^{sat}(T_0)$

On obtient 0.2ml et 1.7ml (0.47*0.51, 0.53*3.3)

Question : Quels sont les échanges d'énergie développés au cours de la transformation suivante. 
I : Vapeur, p_i = 18 bar, T_i= 18°
Transformation isotherme qs
L+V, p_f = p_{sat}(18°0), V_F = V0, T_F = 18°

Ce sont deux états quelconques.

Sur le chemin, on est isotherme mais on a pas $\Delta U=0$ car pas GP.


\end{document}


